\section{Introduction}

\subsection{Classification under equivalence: a millennial question}
\label{ssec:classification}

The question of which discrete objects fit canonically inside which
continuum, and how the continuum organises them by its symmetry, is
older than its modern formalisation.  Vitruvius (\emph{De
Architectura}, c.~30~BC) inscribes the human body simultaneously
within a square and a circle sharing a common centre: the body itself
plays the role of \emph{modulus}, a system of an origin, a radius, an
angle, and a magnitude, mediating between the square as discrete
arithmetic substrate (four elements, four cardinal directions, the
binary doubling $1\mapsto 2\mapsto 4$) and the circle as continuous
rotational symmetry.  Cesare Cesariano's 1521 commentary on Vitruvius
records the same module geometrically through a sequence of nested
regular polygons sharing axis (the origin) and base vertex (the
radius at angle zero), from triangle to dodecagon: each polygon
discretises the circle at the angle $2\pi/n$, all fitting inside the
same continuous rotational ambient.  The polygonal sequence is a
geometric record of a structural question: given an origin, a radius,
and a magnitude, which angular discretisations admit canonical
inscription inside the circle's rotational symmetry, and in what form.

Arguably the evolution of the same ancient structural question is
taken up in two foundational modern treatises.  Euler's
\emph{Introductio in analysin infinitorum} (1748)---characterised by
Boyer at the 1950 ICM as ``the foremost textbook of modern
times''~\cite{boyer-1951}---establishes the function as the central
object of analysis and gives a systematic analytic treatment of the
exponential and trigonometric identities, including
$e^{i\theta}=\cos\theta+i\sin\theta$---an analytic rotor of origin~$0$,
radius~$1$, angle~$\theta$, and unit magnitude---whose inverse form
$i\theta=\log(\cos\theta+i\sin\theta)$ had been known to Cotes
in~1714.  Gauss's \emph{Disquisitiones Arithmeticae} (1801, drafted
from~1796) gives the systematic arithmetic counterpart, consolidating
the work of Fermat, Euler, Lagrange, and Legendre into a unified
framework---including in particular the arithmetic of the cyclotomic
roots $\zeta_{n}=e^{2\pi i/n}$, the angles $2\pi/n$ of the regular
$n$-gon treated as algebraic integers---that, in Kline's phrasing,
``began the modern theory of numbers''~\cite{kline-1972} and set the
directions of subsequent work in the subject.

\subsection{Euler 1748: the analytic-geometric substrate}
\label{ssec:euler-substrate}

The identity
\begin{equation}\label{eq:euler}
  e^{i\theta} = \cos\theta + i\sin\theta
\end{equation}
defines the one-dimensional complex rotor: cosine and sine are the
Cartesian projections of $e^{i\theta}$ onto the real and imaginary
axes.  The
unit circle $S^{1}\subset\mathbb{C}$ is the locus of $e^{i\theta}$, and
$2\pi$ is the length of one revolution.  Within $S^{1}$, the
$n$-th roots of unity $\zeta_{n} = e^{2\pi i/n}$ are the vertices of
the regular $n$-gon: the polygonal discretisation of the circle
\emph{is} the theory of roots of unity.  Cesariano's nested polygons
are the geometric trace of what \eqref{eq:euler} encodes analytically.

The identity $e^{i\pi}+1=0$ is the evaluation of \eqref{eq:euler} at
$\theta=\pi$: an exact identity in $\mathbb{C}$ in which several
independently defined objects---the analytic exponential, the
algebraic unit~$i$, the geometric constant~$\pi$, and the arithmetic
units $1$ and~$0$---coincide simultaneously, in a way not implied by
their definitions taken individually.  The content of~\eqref{eq:euler}
is this simultaneous coincidence of independent identities at a
single object.

Subsequent mathematical work over the analytic-geometric substrate
of~\eqref{eq:euler} has produced further simultaneous identities and
related results of varying form (reciprocity laws, congruences,
equivalences, motivic realisations).  Immediately after Euler, Gauss (\emph{Disquisitiones
Arithmeticae}, 1801) established structural identities connecting the
roots of unity defined by~\eqref{eq:euler} to arithmetic data modulo
primes, opening the systematic arithmetic exploitation of the
substrate.  Over the two centuries that separate Gauss from the
present, this line has multiplied through successive
nineteenth- and twentieth-century developments; an exhaustive
account is beyond the scope of this note.  Several examples
illustrate the range.  In quantum mechanics, the Feynman path
integral expresses the propagation amplitude as a sum
$\int e^{iS[\gamma]/\hbar}\mathcal{D}\gamma$ in which the Eulerian
exponential of the action governs the analytic structure of the
theory.  In harmonic analysis on locally compact abelian groups
(Pontryagin), the characters
$\chi_\xi(x)=e^{2\pi i\langle x,\xi\rangle}$ realise the duality
between a group and its dual; the parallel theory of unitary
representations of compact Lie groups proceeds through characters
valued in the unit circle.  The Hilbert--Pólya conjecture proposes
that the nontrivial zeros of $\zeta(s)$ are the eigenvalues of a
self-adjoint operator $H$, whose associated unitary group
$\{e^{itH}\}_{t\in\mathbb{R}}$ would realise the spectrum as a
one-parameter rotation; the Selberg trace formula relates the
spectrum of the Laplacian on a hyperbolic surface to its closed
geodesics through analogous exponential sums.  In arithmetic
geometry, the Weil cohomologies place the roots of unity
$\zeta_n=e^{2\pi i/n}$ at the centre of the Frobenius action on
étale cohomology, and Grothendieck's theory of motives organises
periods $\int_\gamma\omega$---including the elementary
$\oint dz/z=2\pi i$---as fundamental invariants.  More recently,
identifications of unification across arithmetic, tropical
geometry, and motivic theory have been pursued over the
$\mathbb{F}_{1}$ substrate
(Connes--Consani, 2014--2021~\cite{connes-consani-site}).

Across these developments, arithmetic content otherwise undecidable in
the sense of Gödel~\cite{godel-1931} becomes accessible through the
analytic-geometric substrate of~\eqref{eq:euler}:
Tarski~\cite{tarski-1948} showed that elementary geometry over the
real-closed field admits a decision procedure, by contrast with
elementary arithmetic.  The field-with-one-element programme initiated
by Tits~\cite{tits-1957} and developed by Manin~\cite{manin-1995} among
others pursues a geometry over a hypothetical ``field with one
element'', under which $\mathbb{Z}$ would behave as a curve and roots of
unity supply the arithmetic with $\mathbb{F}_{1}$-structure, so as to
carry methods native to geometry over finite fields into the study
of~$\mathbb{Z}$.

This aim descends from Weil's ``Rosetta Stone''.  In a 1940 letter to his
sister Simone~\cite{weil-1940}, Weil set out an analogy among three
domains---number fields, function fields over a finite field, and Riemann
surfaces---read as three columns of a single text.  In the middle column he
was able to prove the Riemann hypothesis in its form for curves over finite
fields~\cite{weil-1948}; the ambition of the field-with-one-element
programme is to endow $\mathbb{Z}$ with enough geometry that this proof, or
an analogue of it, may be transported to the number-field column and settle
the classical Riemann hypothesis.  The Rosetta Stone is itself an instance
of the classification under equivalence of~\Cref{ssec:classification}: one
arithmetic content read across representations that are not identical but
structurally the same, the decisive work lying in the translation---the
equivalences---between them.

The present note operates from a particular angle on this substrate.
In some sense closer to the modulus of \Cref{ssec:classification} than
to those modern
extensions, it uses the golden ratio $\golden=(1+\sqrt{5})/2$
together with the rotor of Euler's identity to articulate identities
between the analytic and the arithmetic through the geometric
structure of~$\golden$, exploring the relations that $\golden$
admits between $e^{i\theta}$, $\mathbb{R}$, and $\mathbb{C}$.

In what follows the golden ratio $\golden$ enters the argument as a
structural operator between discrete polygonal geometry---specifically
the pentagon and decagon---and the continuous rotational symmetry of
the unit circle.  The relations exploited here are inspired by the
proportional compass of classical artistic practice, their use of
ratios based on the human body and the much later invention of the
\emph{sectio aurea} tool: not as an artisan's empirical
approximation, but as an exact coupling between geometric magnitudes
and algebraic roots.  Grisales Herrera~\cite{grisales} establishes
one face of this coupling between $\golden$ and the integers via the
eight reduced residue classes of $\Ztwenty$; a parallel
angular-rotation framework for primes is developed by Acevedo
Agudelo~\cite{acevedo-2021}.  The present note records, in the
analytic formulation of Euler's identity~\eqref{eq:euler}, the
broader configuration of which these facts are part.

For the golden ratio, the starting point is geometric: among the
vertices $\{\zeta_{10}^{k}\}_{k=0}^{9}$ of the regular decagon
defined by~\eqref{eq:euler} at $n=10$, the sum of two conjugate
vertices equals~$\golden$,
\[
  \golden \;=\; \zeta_{10}+\zeta_{10}^{-1}, \qquad \zeta_{10}=e^{i\pi/5}
  \qquad\text{(clause (i) of \Cref{thm:coincidence-intro})}.
\]
This inscription places $\golden$ inside the structure
of~\eqref{eq:euler} as a canonical point.  Three further inscriptions
follow from this one and constitute the content of the present note.
First, the real part of the same vertex sum gives the trigonometric
identity $\golden = 2\cos(\pi/5)$ in~$\mathbb{R}$ (clause~(iii)).
Second, the multiplicative line $\mathbb{R}_{>0}$ admits a
transcendental period $\lambda_{\log}=\log_{\golden}2$ that plays the
structural role of the $2\pi$ of~\eqref{eq:euler} and closes a turn
between the algebraic objects~$\golden$ and~$2$,
\[
  \golden^{\lambda_{\log}} \;=\; 2 \qquad\text{(clause (iv))},
\]
a real-multiplicative companion to~\eqref{eq:euler}.  Third, the
cyclotomic field $\mathbb{Q}(\zeta_{20})$ where $\golden$ resides has
Galois group $\Ztwenty$ of order~$8$ (clause~(ix)), and the Mersenne
primes $\Mp$ distribute over the structural fraction
\[
  M_{2}/|\Ztwenty| \;=\; 3/8 \qquad\text{(clause (xi))}
\]
of those eight Galois classes.  These five inscriptions, together
with seven further identities, inclusions, and identifications in the
same canonical structures (clauses (ii), (v), (vi), (vii), (viii),
(x), (xii) of \Cref{thm:coincidence-intro}), are collected in the
single formal statement~\Cref{thm:coincidence-intro} below and
verified in Lean~4.  The relevance of $\golden$ at this point in the
argument is that the inscription $\golden=\zeta_{10}+\zeta_{10}^{-1}$
in~\eqref{eq:euler} propagates---trigonometrically into~$\mathbb{R}$,
multiplicatively into~$\mathbb{R}_{>0}$, and arithmetically into the
Galois group of the cyclotomic extension---producing twelve coincident
identities under a single algebraic object.

\subsection{A simultaneous identities theorem for
\texorpdfstring{$\golden$}{phi}}\label{ssec:coincidence-intro}

By a \emph{simultaneous identities theorem} we mean a statement
asserting that a mathematical object satisfies simultaneously a
finite collection of identities provable in distinct canonical
structures, with the simultaneous validity of the collection as the
content of the statement.  The corresponding statement
for~$\golden$ is the following \Cref{thm:coincidence-intro}.  Its
twelve clauses are proved in \Cref{sec:universal}
and~\Cref{sec:triple} (clauses (i)--(iii),~(ix)
in~\Cref{rem:pentagon}, clause~(iv) in~\Cref{lem:binary}, clause~(v)
in~\Cref{thm:universal}, clause~(vi)
in~\Cref{thm:concentration-main}, clause~(vii)
in~\Cref{thm:three-main}); the four remaining clauses (viii), (x),
(xi), (xii) have short proofs but each registers a structural
knotting articulated after the statement.

The individual proofs of the clauses are short---several require only
definitional manipulation or elementary modular computation.  The
substance of the statement does not reside in the difficulty of the
separate proofs but in the simultaneity of their validity under a
single~$\golden$, and in the absence of a prior reference recording
them jointly.  The situation is analogous to that of Euler's identity
$e^{i\pi}+1=0$, whose modern proof is the evaluation of a series but
whose enunciation revealed a knotting of analysis, algebra, geometry,
and arithmetic not visible in the prior literature.  The formal
verification in Lean~4 secures the simultaneity as a single theorem,
eliminating the possibility that the clauses hold separately but not
jointly.  Throughout, $\iota:\mathbb{R}\hookrightarrow\mathbb{C}$
denotes the canonical embedding and $\golden_\mathbb{C}:=\iota(\golden)$
its image in~$\mathbb{C}$ under this embedding.

\begin{Theorem}[Simultaneous identities for $\golden$ (formal version:
\Cref{thm:euler-connector})]\label{thm:coincidence-intro}
A single $\golden$ satisfies simultaneously the following twelve
identities, inclusions, and identifications, organised by the
canonical structure they inhabit:

\smallskip
\noindent\emph{In the complex rotor of Euler's identity
$\{e^{i\theta}\}\subset\mathbb{C}$:}
\begin{enumerate}[topsep=0.2em,itemsep=0.1em,leftmargin=2.3em]
  \item[\textup{(i)}] \emph{Cyclotomic inscription:}
    $\golden_\mathbb{C} = \zeta_{10}+\zeta_{10}^{-1}$, with
    $\zeta_{10}:=e^{i\pi/5}$.
  \item[\textup{(ii)}] \emph{Tenth-root closure:} $\zeta_{10}^{10}=1$.
  \item[\textup{(xii)}] \emph{Rotor discretisation:} for every
    $k\in\mathbb{N}$, $(\zeta_{10}^{k})^{10}=1$.
\end{enumerate}

\smallskip
\noindent\emph{In the real trigonometric line, with cross-identification
to the complex inscription:}
\begin{enumerate}[topsep=0.2em,itemsep=0.1em,leftmargin=2.3em]
  \item[\textup{(iii)}] \emph{Pentagonal trigonometric identity:}
    $\golden = 2\cos(\pi/5)$.
  \item[\textup{(viii)}] \emph{Real-complex identification:}
    $\iota(\golden) = \golden_\mathbb{C}$.
\end{enumerate}

\smallskip
\noindent\emph{In the real-multiplicative line $\mathbb{R}_{>0}$:}
\begin{enumerate}[topsep=0.2em,itemsep=0.1em,leftmargin=2.3em]
  \item[\textup{(iv)}] \emph{Logarithmic binary bridge:}
    $\golden^{\lambda_{\log}}=2$, where
    $\lambda_{\log}:=\log_{\golden}2$.
  \item[\textup{(v)}] \emph{Universal real-multiplicative identity:}
    $\forall p\in\mathbb{R},\;3\,\golden^{\sigma(p)}=2^{p}$, where
    $\sigma(p):=p\,\lambda_{\log}-\log_{\golden}3$.
\end{enumerate}

\smallskip
\noindent\emph{In the arithmetic of Mersenne modulo 20:}
\begin{enumerate}[topsep=0.2em,itemsep=0.1em,leftmargin=2.3em]
  \item[\textup{(vi)}] \emph{Modular concentration:} for every prime
    $p$, $\Mp\bmod 20\in\{3,7,11\}$.
  \item[\textup{(vii)}] \emph{Triple ternary identification:}
    $3 = M_{2} = |\{3,7,11\}|$.
\end{enumerate}

\smallskip
\noindent\emph{In the Galois group of the cyclotomic extension
$\mathbb{Q}(\zeta_{20})/\mathbb{Q}$, with quantitative ratio to the
arithmetic locus:}
\begin{enumerate}[topsep=0.2em,itemsep=0.1em,leftmargin=2.3em]
  \item[\textup{(ix)}] \emph{Galois order:}
    $|\Ztwenty|=8$.
  \item[\textup{(x)}] \emph{Galois sub-object inclusion:}
    $\{3,7,11\}\subseteq\Ztwenty$.
  \item[\textup{(xi)}] \emph{Quantitative ratio:}
    $|\{3,7,11\}|/|\Ztwenty|=M_{2}/8=3/8$, equivalently
    $|\{3,7,11\}|\cdot 8 = M_{2}\cdot|\Ztwenty|$.
\end{enumerate}
\end{Theorem}

No proper sub-collection of the twelve clauses constitutes the content;
the content is the simultaneous holding of all twelve under a single
$\golden$ and in relation to Euler's identity~\eqref{eq:euler}, where
the constants $\zeta_{10}=e^{i\pi/5}$, $\lambda_{\log}=\log_{\golden}2$,
and $\sigma(p)=p\,\lambda_{\log}-\log_{\golden}3$ are formally introduced
in~\Cref{sec:universal} (\Cref{def:sigma}) and~\Cref{rem:pentagon}.
The five groupings above record the five canonical structures in
which $\golden$ is independently inscribed: the complex rotor with
its tenth-root discretisation, the real trigonometric line, the
real-multiplicative line, the arithmetic of Mersenne modulo~20, and
the Galois group of the cyclotomic decagonal extension.  The
cross-identifications (clauses (viii) and (xi)) record that the
inscriptions are inscriptions of \emph{the same}~$\golden$ across
distinct substrates and that the arithmetic locus sits as a
quantitatively specific sub-object of the Galois group.  \Cref{ssec:characterization} makes this precise as a
consistency theorem: the factor~$3$, the base~$\golden$, and the slope
$\lambda_{\log}$ are each pinned by a theorem external to the central
identity, and the identity is the unique point at which they agree.

\begin{figure}[H]
  \centering
  \includegraphics[width=\textwidth]{cd_correspondance.pdf}
  \caption{The twelve simultaneous identities of~\Cref{thm:coincidence-intro}
  (clauses (i)--(xii)), together with the reciprocal turn
  $\golden^{-\lambda_{\log}}=\tfrac12$ of~\Cref{sec:spectral}, as a single
  commutative diagram.  Every directed path commutes to the universal
  identity $3\,\golden^{\sigma(p)}=2^{p}$ (green) and to the reciprocal turn
  $\tfrac12=\golden^{-\lambda_{\log}}$ (red, left); the factor $3=M_{2}$ and
  the half each appear as a convergence of several identifications under the
  single algebraic object~$\golden$.}
  \label{fig:coincidence-diagram}
\end{figure}

Clauses (vi), (vii), (x), (xi), taken in isolation, fall in the
genealogy of arithmetic identities over
$\mathbb{Q}(\sqrt{5})$ companion to Lucas's 1878 Fibonacci
congruence and to Grisales Herrera's classification of the eight
residue classes of $\Ztwenty$ into split and inert types
in~$\mathbb{Q}(\sqrt{5})$~\cite{grisales}.  They enter
\Cref{thm:coincidence-intro} as four of the twelve simultaneous clauses
of the simultaneous identities theorem, not as the global content of the
theorem.  The splitting type of each prime $\Mp$ in $\mathbb{Q}(\sqrt5)$ is
developed in~\Cref{sec:fibonacci} via Lucas's congruence, where it is
shown to be governed by the exponent $p\bmod4$ and to single out the
inert primes as those satisfying the Lucas--Lehmer divisibility
$\Mp\mid F_{\Mp+1}$.

Beyond these twelve simultaneous clauses, the same multiplicative period
$\lambda_{\log}$, read in reverse, closes a reciprocal turn
$\golden^{-\lambda_{\log}}=\tfrac12$.  This value is not a thirteenth
clause but a point already inscribed in the identity---the reciprocal of
clause~(iv), the $p=-1$ node of the universal identity~(v), and the
multiplicative-line image of the rotor's conjugate root (clause~(i))---%
developed in~\Cref{sec:spectral}, where it is also identified as the
value at which the companion operator paper anchors its Hilbert--P\'olya
operator.

All statements of this note are formally verified in
\texttt{golden\_mersenne.lean}, a self-contained Lean~4
file built against Mathlib4 (toolchain~\texttt{v4.28.0})~\cite{mathlib};
no \texttt{sorry} occurs in any proof.  Two external inputs enter as
axioms: the Gelfond--Schneider theorem (not yet available in Mathlib
mainline), which underlies the transcendence of $\lambda_{\log}$, and the
44 Lucas--Lehmer~\cite{lehmer} primalities $M_{p}\;\text{prime}$ for
$61\leq p\leq 136\,279\,841$ (GIMPS~\cite{gimps}).  The relation
$\golden^{2}=\golden+1$ is a theorem, not an axiom.

